\subsection{The parameter inequalities}
\label{sec:param-lists}

The common setting in Section~\ref{sec:relay} uses the choices (P0)--(P6)
below. They are made in the order of Table~\ref{tab:param}. The inequalities
ensure the hypotheses of the estimates used in the relay.

\begin{definition}
\label{def:P0-P4}
\emergencystretch=3em
The \textbf{structural list} consists of the following ordered choices.
The auxiliary constants are defined at the cited locations.
\begin{description}[font=\normalfont]
\item[\textup{(P0.1)}] $\Phi$, $a$, $\Gone$ of Definition~I.\ref{I-def:ring}; $\lamref$, $\sigma$,
  $\sigb$ of Definition~I.\ref{I-def:eigendata}.
\item[\textup{(P0.2)}] An integer $s$ and a number $\sigma'$ as in (SA1); then
  $c_\gamma$ of Appendix~\ref{sec:relay-constants}.
\item[\textup{(P0.3)}] The integers $N_G$, $M$, the constants
  $\dH,C_2,C_3$ and $\sigma_{\max}$, and the functions
  $\kappa_0(\cdot)$, $\bar t_1(\cdot)$, $C(\cdot)$ of Theorem~\ref{thm:active};
  $\kL$ and $\Cseed$ of Lemma~\ref{lem:seed}; $\|\ellst\|$, $\Cls$ and $\Cp$ of
  Lemma~I.\ref{I-lem:uniform}.
\item[\textup{(P0.4)}] The integers $s_h$, $N_Q$, $M_h$ and the constants
  $\majst,C_Q$ of (SA2).
\item[\textup{(P0.5)}] The common velocity constant $\Cvel$ of
  Lemmas~\ref{lem:transfer}, \ref{lem:inner-field} and~\ref{lem:field-bound},
  at order $M_h$; the absolute constant $\Cstar$ of Lemma~\ref{lem:speed}; and
  $\Cmom$ of Proposition~\ref{prop:confine}.
\item[\textup{(P0.6)}] $c_N$, $C_N$, $\cout$ of Lemma~I.\ref{I-lem:uniform};
  the constants of Lemma~\ref{lem:kicktrue} are chosen at (P4).
\item[\textup{(P1)}] Choose $\delta_0\in(0,\dH]$ sufficiently small that
  $\frac{16}{15}\cdot\frac{2\pi C_3\delta_0^2}{\sigma}\le\frac14\cdot\frac{15}{16}\cdot\frac{\cout\delta_0}{20|\lamref|}$
  and $\Cp C_2\delta_0\le\frac{\cout}{16\pi}$.
\item[\textup{(P2)}] $\mu\in(1,\mu_{\max}]$, with
  $\mu_{\max}=1+\min\big(c'\delta_0,\frac1{16}\big)$ and
  a fixed $c'>0$ sufficiently small that
  $c'\le\frac{15}{16}\cdot\frac{\cout}{10|\lamref|\,\Cmom\,\Ip{3}{\sigma/4}}$; then
  $\vartheta_\mu$, $\gamma_\mu$, $\gamma'_\mu$ of Appendix~\ref{sec:relay-constants},
  and $\CDuh$, $\KO$ of Lemma~\ref{lem:O}.
\item[\textup{(P3)}] $\Tpin>0$ with $\rT(\Tpin)\le\frac1{64}$, with the function
  $\rT$ of Definition~\ref{def:standing}. Then $C_D(\Tpin)$
  (Proposition~\ref{prop:relay}); $K_1$, $\Abar$, $\Klo$, $\Kup$
  (Definition~\ref{def:K1});
  $\Tz$ (Lemma~\ref{lem:time}) and
  choose $C_\tau\ge\max\{3,\frac{15}{\sigma'}\}$.
\item[\textup{(P4)}] Choose $C_0$ so that
  \[
  \begin{gathered}
  C_0\ge\max\{1,2\Cseed,K_1C_Q\},\\
  C_0\ge\Cvel \,\mu\,(2+2\Cstar\Tpin)^3(1+2\Cstar)^3/\Klo.
  \end{gathered}
  \]
  Then fix $\kappa_0(C_0)$, $\bar t_1(C_0)$ and $C(C_0)$ of
  Theorem~\ref{thm:active}, the common constant $\Crelay(C_0,K_1)$ of
  Proposition~\ref{prop:relay} and Lemma~\ref{lem:Oprime}, and $\Cerr$ of
  Lemma~\ref{lem:kicktrue}.
\end{description}
\end{definition}

The constant $C_0$ serves (Init) and (Field) in Theorem~\ref{thm:active}.
The seed estimate bounds the initial perturbation by $2\Cseed\ksh_n$.
The hand-off estimate in Appendix~\ref{sec:bootstrap-handoff-checks}
bounds it by $K_1C_Q\ksh_k$ at ring $k$'s activation. The inner-field bound
in Appendix~\ref{sec:bootstrap-active-checks} is the coefficient in the
second inequality of (P4) times $\ksh_k\jap{A_k(t-\tsh_k)}^3$. These are
the bounds required in the local coordinate $u=A_k(t-\tsh_k)$.
The coefficients that $C_0$ must dominate are fixed at (P0)--(P3).
The constants $\Crelay$ and $\Cerr$ are fixed after $C_0$ at (P4), before
choosing the seed and radius ratios.

\begin{definition}
\label{def:P5}
The \textbf{$\kappa$-list (P5)}: choose $\kseed>0$ so small that, for every
$\kappa\in(0,2\kseed]$, the following hold.
\begin{description}[font=\normalfont]
\item[\textup{($\kappa$1)}] $\kappa\le\delta_04^{-\mu/(\mu-1)}$.
\item[\textup{($\kappa$2)}] \begin{multline*}
  \frac{2\pi C_3\kappa^{1/2}\mu}{\frac34\sigma(\mu-1)}+\Cerr\Big(\frac\kappa{\delta_0}\Big)^{\gamma_\mu}\log^4\frac{\delta_0}\kappa\\
  +\Cerr\big[1+\rT(\Tpin)\Ip{3}{\tfrac\sigma8}\big]\Big(\frac\kappa{\delta_0}\Big)^{\vartheta_\mu}\log^3\frac{\delta_0}\kappa
  \le\frac12\cdot\frac{15}{16}\cdot\frac{\cout}{20|\lamref|(\mu-1)}.
  \end{multline*}
\item[\textup{($\kappa$3)}] $\Cerr(\kappa/\delta_0)^{\vartheta_\mu}\le\frac{15}{16}\cdot\frac{\cout}{20|\lamref|(\mu-1)}$.
\item[\textup{($\kappa$4)}] $\kappa\le\kappa_0(C_0)$.
\item[\textup{($\kappa$5)}] For every $G\in[0,\Gone]$, with $\hat T=\sigma(G)^{-1}\log(\delta_0/\kappa)$
  and $\rhohat$ the transient of Theorem~\ref{thm:active} for $(C_0,G)$:
  $\Cp C_2\big(\rhohat(\hat T)+\kappa^{1/2}\big)\le\frac{\cout}{16\pi}$.
\item[\textup{($\kappa$6)}] $\kappa\le\kL$.
\item[\textup{($\kappa$7)}] $\Crelay(C_0,K_1)\big(1+\frac1{\sigb(\mu-1)}\big)(\kappa/\delta_0)^{\gamma'_\mu}\langle\log(\delta_0/\kappa)\rangle^4\le\frac1{\sigb(\mu-1)}$.
\end{description}
\end{definition}

In (P5), $\kappa$ stands for an amplitude $\ksh_k$ with which a ring starts its
active phase and, in ($\kappa$1) and ($\kappa$6), also for $2\kseed$.
Theorem~\ref{thm:active} is applied to ring $k$ with $\kappa=\ksh_k$, and
Lemma~\ref{lem:kicktrue}, at level $k$, with $\kappa=\ksh_{k+1}$.
Conditions ($\kappa$1)--($\kappa$3) are hypotheses (c)--(e) of
Lemma~\ref{lem:kicktrue}, ($\kappa$4) is hypothesis (Small) of
Theorem~\ref{thm:active}, ($\kappa$5) is used in the proof of
Proposition~\ref{prop:reduction}, and ($\kappa$6) makes the seed of
Definition~\ref{def:approx} admissible. The seed estimate
\eqref{eq:boot-start} and the hand-off estimates of
Appendix~\ref{sec:bootstrap-handoff-checks} put every starting amplitude
in $(0,2\kseed]$, so the conditions are imposed on that whole interval.

The formula for $K_1$ in Definition~\ref{def:K1} uses quantities fixed at
(P0)--(P3). The $+1$ in $4C_2+1$ allows the remainder $\Xi_{j+1}$ in
Proposition~\ref{prop:relay} to be as large as $\delta_0/(\sigb(\mu-1))$.
After $C_0$ and $\Crelay(C_0,K_1)$ are fixed, ($\kappa$7) ensures that this
allowance suffices for every $\kappa=\ksh_{j+1}\in(0,2\kseed]$.
The comparison of exponents is checked in
Appendix~\ref{sec:bootstrap-handoff-checks}. The same $K_1$ and $\Abar$
therefore work at every hand-off.

The number $\sbar$, which fixes the ratios $s_k$ of consecutive radii, is chosen
last, because the conditions (B1)--(B11) on it involve constants fixed at all the
earlier steps; each of them holds for every
sufficiently small $\sbar$ (proof of Lemma~\ref{lem:nonempty}).

\begin{definition}
\label{def:P6}
The \textbf{$\sbar$-list (P6)}: choose $\sbar>0$ such that the following hold,
where $s_k=\sbar(2\mu)^{-k}$, $\eta_{b,1}=0$ and $\eta_{b,k}$, $\eta_b$ are as in
Equation~\eqref{eq:etab}.
\begin{description}[font=\normalfont]
\item[\textup{(B1)}] $\sbar\le a/16$ and $\sbar\le\frac12$.
\item[\textup{(B2)}] $\mu^4\sbar^2\le\frac18$.
\item[\textup{(B3)}] $\mu\sbar^2\le\Gone$.
\item[\textup{(B4)}] $\sup_{0\le G\le \mu\sbar^2}|\lamst(G)-\lamref|\le\sigma(\mu-1)/(8\mu)$
  (this is (Det)).
\item[\textup{(B5)}] For every $k\ge1$: $\Lconf_ks_k\le a/8$ and
  $\frac83(\frac2a)^3(\Lconf_k)^2s_k\le\mu\Cstar$, where
  \[
  \Lconf_k=2+2\Cstar\mu\Big[\Tpin+1+\frac1{\sigma'}\Big(5\log\frac1{s_k}+\log_+\frac{\delta_0}{\Klo}\Big)\Big].
  \]
\item[\textup{(B6)}] $\Kup\sbar^5\le\kseed$.
\item[\textup{(B7)}] $16\Cvel C_Q\CDuh\KO\sbar^6\le1$, $\KO\sbar^6\le\frac14$,
  $\KO\mu\sbar^5\le\frac12$, and
  \begin{equation}
  \label{eq:E22}
  \KO\sbar^6\Big[\Tz+2\log\frac1{\sbar}+\frac1{\mu-1}\Big]\le\frac{\mu-1}{4\Cvel C_QC_\tau}\min\Big\{\frac12,\frac1{4\Cls K_1}\Big\}.
  \end{equation}
\item[\textup{(B8)}] $\eta_b\le\frac12$; for every $k\ge2$,
  \[
  \frac{\Cls\eta_{b,k}}{\sigma'}\Big[5\log\frac1{s_k}+\log_+\frac{\delta_0}{\Klo}\Big]\le\frac1{32};
  \]
  and $\Cls\eta_b\sigma'^{-1}\log(2\delta_0/\kseed)\le\frac1{32}$.
\item[\textup{(B9)}] $2\Abar\mu\sbar^5\le\majst$ and $\Cvel (\mu\sbar^2+\sbar)\le1$.
\item[\textup{(B10)}] $4\Cvel \big(\frac{2C_Q}{\sigb(\mu-1)}+1\big)\mu\sbar^5\le\frac12$.
\item[\textup{(B11)}] $\|\ellst\|C_QK_1\Abar\mu\sbar^5\le\frac12$, and
  \[
  \Cerr\eta_b\delta_0\frac{\mu^2}{\sigma^2(\mu-1)^2}+2\pi C_3\eta_b\frac{4\mu\delta_0}{3\sigma(\mu-1)}
  \le\frac14\cdot\frac{15}{16}\cdot\frac{\cout\delta_0}{20|\lamref|(\mu-1)} .
  \]
\end{description}
\end{definition}

The second part of (B1), and (B2), hold automatically (see the remark after
Equation~\eqref{eq:sk-small}); they are listed because Proposition~\ref{prop:confine} and
Lemma~\ref{lem:field-bound} state them as hypotheses.

\subsection{The accumulated outer field}
\label{sec:param-feedback}

The last inequality of (B7), Equation~\eqref{eq:E22}, is used only through the
following lemma. Recall that $\theta_j$ is the factor of Lemma~\ref{lem:O} that
bounds the field of the outer rings on ring $j$ by the Duhamel bound
$\maj_j$, and that $\bar\tau^\sharp_j$ of Lemma~\ref{lem:quiet-time} bounds
physical times in its quiet-ring estimate by $A_jt\le\bar\tau^\sharp_j$.
Thus $\theta_jA_jt\le\theta_j\bar\tau^\sharp_j$. The lemma bounds
this dimensionless product, which must be at most $\frac12$ for hypothesis \textup{(H2)} of
Proposition~\ref{prop:relay} and for the corresponding hypothesis of
Lemma~\ref{lem:Oprime}, and at most $1/(4\Cls K_1)$ for hypothesis
\textup{(ii)} of Lemma~\ref{lem:handoff}.

\begin{lemma}
\label{lem:theta-quiet}
Assume (SA) and \textup{(B7)}, and let $K_1$ be as in Definition~\ref{def:K1},
$\Tz$ and $C_\tau$ as in \textup{(P3)}, $\theta_j$ as in Lemma~\ref{lem:O} and
$\bar\tau^\sharp_j$ as in Lemma~\ref{lem:quiet-time}. Then $\theta_j\bar\tau^\sharp_j\le\min\{\frac12,\frac1{4\Cls K_1}\}$
for every level $1\le j<n$.
\end{lemma}

\begin{proof}
For $j=1$, $\theta_1=0$ by Lemma~\ref{lem:O}, and the claim holds. Let
$2\le j<n$. By Lemma~\ref{lem:O}, $\theta_j=4\Cvel C_Q\KO s_{j-1}^6$.

We first bound $\bar\tau^\sharp_j$ by a multiple of $C_\tau$. By (SA3),
$1<\mu\le\frac{17}{16}$, so $\mu\log(2\mu)\le\frac{17}{16}\log\frac{17}8<1$.
The three terms of $\bar\tau^\sharp_j$ (Lemma~\ref{lem:quiet-time}) are those in
$\Tz$, in $\log(1/s_j)$ and in $(\mu-1)^{-2}$; we compare them with
those of $C_\tau[\Tz+\log(1/s_j)+(\mu-1)^{-1}]/(\mu-1)$. By the definition of
$C_\tau$ in \textup{(P3)} and the bound on $\mu\log(2\mu)$, the coefficients
satisfy $3\le C_\tau$,
$11/\sigma'\le C_\tau$ and $5\mu\log(2\mu)/\sigma'\le5/\sigma'\le C_\tau$;
this gives $\bar\tau^\sharp_j\le C_\tau\big[\Tz+\log(1/s_j)+(\mu-1)^{-1}\big]/(\mu-1)$.
Next we pass from $s_j$ to $s_{j-1}$. Since
$s_{j-1}\le\sbar\le a/16<\frac1{2\mu}$, we have
$\log\frac1{s_j}=\log\frac1{s_{j-1}}+\log(2\mu)\le2\log\frac1{s_{j-1}}$, so
\[
\theta_j\bar\tau^\sharp_j\le\frac{4\Cvel C_QC_\tau}{\mu-1}\,\KO\,g(s_{j-1}),\qquad
g(x)=x^6\Big[\Tz+2\log\frac1x+\frac1{\mu-1}\Big].
\]
It remains to replace $s_{j-1}$ by $\sbar$, since Equation~\eqref{eq:E22}
bounds $\KO g(\sbar)$; for this we show that $g$ is nondecreasing near $0$. For $0<x<1$,
$g'(x)=x^5\big[6\big(\Tz+2\log\frac1x+\frac1{\mu-1}\big)-2\big]$, which is
nonnegative when $\log(1/x)\ge1/6$, because $\Tz>0$ and $\mu>1$. So $g$ is
nondecreasing on $(0,e^{-1/6}]$. Since $s_{j-1}\le\sbar\le a/16<\frac1{16}<e^{-1/6}$,
$g(s_{j-1})\le g(\sbar)$, and Equation~\eqref{eq:E22} gives
\[
\theta_j\bar\tau^\sharp_j\le\frac{4\Cvel C_QC_\tau}{\mu-1}\KO g(\sbar)\le\min\{\frac12,\frac1{4\Cls K_1}\}.\qedhere
\]
\end{proof}

\subsection{Verification of the parameter choices}
\label{sec:param-verification}

\begin{proof}[Detailed verification of Lemma~\ref{lem:nonempty}]
We take the steps in order. At each step the quantities of the earlier steps are
fixed, and we show that the conditions of the step can be met.

\emph{(P0)--(P4).} The constants of (P0) exist and are finite by the results
cited there; in particular, by the formulas of (SA2) and $M\ge s+4$
(Theorem~\ref{thm:active}), $s_h=M-1\ge s+3\ge3$ and $N_Q\sigb\ge\Lam(s_h)$, as
Theorem~\ref{thm:quiet} requires. The interval $\big(\frac9{200},\min(\sigb,\frac34\sigma)\big)$ for $\sigma'$
in (SA1) is nonempty: by Proposition~I.\ref{I-prop:U}\textup{(a)}, $\sigb$ is attained and
$\sigb>\frac{211}{2500}>\frac9{200}$; and
$\frac34\sigma\ge\frac34\sigb>\frac{633}{10000}>\frac9{200}$. Moreover
$c_\gamma>0$, because $\sigma'<\frac34\sigma<\frac{15}{16}\sigma$. The conditions in (P1) hold for sufficiently small $\delta_0>0$ because
$\sigma\ge\sigb>0$, $\cout>0$ and the other constants have already been fixed.
In (P2), $c'>0$ because $\Cmom>0$ and $0<\Ip{3}{\sigma/4}<\infty$, so
$\mu_{\max}>1$. In (P3), $\Tpin$ exists because $\rT(T)=(1+T)^3e^{-\sigma T/4}$
(Definition~\ref{def:standing}) tends to $0$ as $T\to\infty$, as $\sigma>0$. The formulas of (P3) and (P4) have positive finite denominators:
$c_N>0$, $\delta_0>0$, $\sigb>0$, $\mu-1>0$ and $\Klo>0$. So all the
constants of (P0)--(P4) are finite and positive, with $C_0\ge1$ as
Theorem~\ref{thm:active} requires.
Each can be chosen in terms of the quantities in its row of Table~\ref{tab:param}, which are fixed at
earlier rows or earlier in the same step.

\emph{(P5).} Since (P5) is imposed for every $\kappa\in(0,2\kseed]$, it suffices
to find $x_0>0$ such that each condition holds for every $\kappa\in(0,x_0]$; then every $\kseed\in(0,x_0/2]$ is admissible. For ($\kappa$1), ($\kappa$4), ($\kappa$6) this holds
because $\delta_04^{-\mu/(\mu-1)}$, $\kappa_0(C_0)$, $\kL$ are positive. For
($\kappa$2), ($\kappa$3), ($\kappa$7) it suffices that each left side tends to
$0$ as $\kappa\to0$. The left sides are finite sums of terms $C\kappa^{1/2}$ and
$C(\kappa/\delta_0)^\gamma\langle\log(\delta_0/\kappa)\rangle^N$ with $C$ fixed at
(P0)--(P4), integers $N\ge0$, and
$\gamma\in\{\gamma_\mu,\vartheta_\mu,\gamma'_\mu\}$, positive because $\mu>1$ and
$c_\gamma>0$; each term tends to $0$. For ($\kappa$5) we need a bound uniform in
$G$. Let $\kappa<\delta_0$ and $G\in[0,\Gone]$, and put $c=\sigb-\sigma'>0$, where
$\sigma'<\sigb$ by (SA1). By Theorem~\ref{thm:active},
$t_1(C_0,G)\le\bar t_1(C_0)$, and $\sigma(G)\in[\sigb,\sigma_{\max}]$, so
$\hat T\ge\sigma_{\max}^{-1}\log(\delta_0/\kappa)$. If moreover
$\kappa\le\delta_0e^{-\sigma_{\max}\bar t_1(C_0)}$, then $\hat T\ge\bar t_1(C_0)$, so
$y=\hat T-t_1(C_0,G)\ge\hat T-\bar t_1(C_0)\ge0$. Hence $\hat T$ lies in the range
$t\ge t_1$ of the second formula for the transient in Theorem~\ref{thm:active},
and, as $\sigma(G)-\sigma'\ge c$,
\[
\rhohat(\hat T)=\jap{y}^3e^{-(\sigma(G)-\sigma')y}\le\jap{y}^3e^{-cy}\le C(c)e^{-cy/2}
\le C(c)e^{c\bar t_1(C_0)/2}\Big(\frac\kappa{\delta_0}\Big)^{c/(2\sigma_{\max})},
\]
where $C(c)$ can be chosen because $\sup_{y\ge0}\jap{y}^3e^{-cy/2}<\infty$; the last step uses
$y\ge\hat T-\bar t_1(C_0)$ and the lower bound on $\hat T$. The right side does not depend on
$G$ and tends to $0$ as $\kappa\to0$, and so does $\kappa^{1/2}$; so ($\kappa$5)
holds for all small $\kappa$, uniformly in $G\in[0,\Gone]$. We take for $x_0$
the least of the thresholds found for ($\kappa$1)--($\kappa$7).

\emph{(P6)}, with (P0)--(P5) fixed. It suffices that each condition holds for
every $\sbar$ in some interval $(0,\sbar_i]$, $\sbar_i>0$; then
every $\sbar\in(0,\min_i\sbar_i]$ is admissible. The conditions (B1)--(B3), (B6), (B7), (B9),
(B10), (B11), the first part of (B8) and its last part (with $\kseed$ fixed)
have left sides that are finite sums of terms $C\sbar^\alpha\log^N(1/\sbar)$ with
$\alpha>0$, $N\in\{0,1\}$ and $C$ fixed at (P0)--(P5); here we use
$\eta_b=\Crelay\sbar^6/(\sigb(\mu-1))$. Each such term tends to $0$ as
$\sbar\to0$, while the right sides are fixed positive numbers. For (B4), the
supremum tends to $0$ as $\sbar\to0$, because $\lamst$ is continuous at $G=0$
(Proposition~I.\ref{I-prop:U}\textup{(a)}) and $\lamst(0)=\lamref$; the right side
$\sigma(\mu-1)/(8\mu)$ is positive. The conditions (B5) and (B8) are imposed at
every level $k$, and we reduce each of them to one condition on $\sbar$. For
(B5), put
$\alpha_0=2+2\Cstar\mu\big[\Tpin+1+\sigma'^{-1}\log_+(\delta_0/\Klo)\big]$ and
$\beta_0=10\Cstar\mu/\sigma'$, so that $\Lconf_k=\alpha_0+\beta_0\log(1/s_k)$.
The functions $x\mapsto x(\alpha_0+\beta_0\log(1/x))$ and
$x\mapsto x(\alpha_0+\beta_0\log(1/x))^2$
tend to $0$ as $x\to0$, so there is $\sbar_{\mathrm{B5}}>0$ such that both inequalities of
(B5) hold at level $k$ whenever $s_k\le\sbar_{\mathrm{B5}}$. Since $s_k\le\sbar$ for every
$k$, (B5) holds for every $\sbar\le\sbar_{\mathrm{B5}}$. For (B8) and $k\ge2$, we have
$s_k=s_{k-1}/(2\mu)$, and, by the formula for $\eta_{b,k}$ in
Equation~\eqref{eq:etab}, the $k$-th inequality reads $F(s_{k-1})\le\frac1{32}$
with
\[
F(x)=\frac{\Cls\Crelay}{\sigma'\sigb(\mu-1)}\,x^6\Big[5\Big(\log\frac1x+\log(2\mu)\Big)+\log_+\frac{\delta_0}{\Klo}\Big].
\]
Since $F(x)\to0$ as $x\to0$, there is $\sbar_{\mathrm{B8}}>0$ with $F\le\frac1{32}$ on
$(0,\sbar_{\mathrm{B8}}]$, and, since $s_{k-1}\le\sbar$ for $k\ge2$, the
inequalities hold for all $k\ge2$ once $\sbar\le\sbar_{\mathrm{B8}}$.

No formula or condition of (P0)--(P6) contains $n$, $A_1$ or $\omega^{(n)}$; the
levels in (B5) and (B8) range over all integers $k\ge1$ and $k\ge2$, not only
up to $n$.

\emph{(SA) and (Det).} Assume (P0)--(P6) and let $n\ge2$ and $A_1>0$. (SA1) holds by (P0.2) and (P0.3), and (SA2) by (P0.4) and (P0.5). For
(SA3): $\mu\in(1,\frac{17}{16}]$ because
$\mu_{\max}\le1+\frac1{16}$; $\sbar\in(0,a/16]$ and $\mu\sbar^2\le\Gone$ by (B1) and (B3);
$\delta_0\in(0,\dH]$ by (P1); $\Tpin>0$ by (P3); and $\kseed\in(0,\kL]$ by
($\kappa$6) at $\kappa=2\kseed$. The objects of
Definitions~\ref{def:datum}--\ref{def:clock} are then defined with these
parameters: Definition~\ref{def:datum} needs $\mu\in(1,2]$ and $\sbar\in(0,a/16]$,
Definition~\ref{def:approx} needs $\kseed\in(0,\kL]$, and
Definition~\ref{def:clock} needs $\delta_0,\Tpin>0$ and $G_k\le\Gone$ for
$1\le k\le n$, which holds by Lemma~\ref{lem:inner-circ} and (B3). Hypothesis
(Det) is (B4).
\end{proof}
